\newpage

\section{\huge \textcolor{blue}{\hspace*{0.2em}极限环及方程组关系}}
\subsection{\LARGE \hspace*{0.1em}极限环}
\subsubsection{\large 极限环的定义}
对如下方程组:
\[
\begin{cases}
    x' = f(x,y) \\
    y' = g(x,y)
\end{cases}
\]
几何: 先建立速度场
\[ \vec{F} = f\hat{i} + g\hat{j} \]
正常有起始点的曲线\\
从解的角度:\\
常数解 $
\begin{cases}
    x(t) = x_0 \\
    y(t) = y_0
\end{cases}
$\\[0.3em]
从场的角度:\\
$ \vec{F} = \vec{0} $ , 场为零处 或 那无速度向量
\[ f(x_0,y_0) = 0, g(x_0,y_0) = 0 \]
现在考虑轨迹的另一种可能"不停止的运动", 称闭合轨迹.\\
\begin{figure}[htbp]
    \centering
\includegraphics[width=0.3\textwidth]{C:/Users/37/Desktop/Notes on Common Calculus Equations/Notes and pictures/cwf13.1.1.png}
    \caption{闭合轨迹}
    \label{fig:cwf13.1.1}
\end{figure}

eg: 表示一个周期性地回到原始状态的方程组:
\[
\begin{cases}
    x' = y \\
    y' = -x
\end{cases} \hspace*{1em} 
\begin{bmatrix}
    0 & 1 \\
    -1 & 0
\end{bmatrix} \hspace*{1em}
\lambda = \pm i
\]
\[
\vec{X} = c_1 
\begin{bmatrix}
    \cos t \\
    -\sin t
\end{bmatrix} + c_2
\begin{bmatrix}
    \sin t \\
    \cos t
\end{bmatrix}
\]
\begin{figure}[htbp]
    \centering
\includegraphics[width=0.4\textwidth]{C:/Users/37/Desktop/Notes on Common Calculus Equations/Notes and pictures/cwf13.1.2.png}
    \caption{周期闭合轨迹}
    \label{fig:cwf13.1.2}
\end{figure}

现可进一步探讨 极限环:\\
极限环 $=$ 闭合轨迹 $+$ 孤立的, 稳定的性质.\\
\begin{figure}[htbp]
    \centering
\includegraphics[width=0.5\textwidth]{C:/Users/37/Desktop/Notes on Common Calculus Equations/Notes and pictures/cwf13.1.3.png}
    \caption{极限环}
    \label{fig:cwf13.1.3}
\end{figure}

图中每条轨迹都不会相交, 但会趋近极限环, 
若存在不稳定点, 曲线为远离极限环而非靠近, 而不稳定点不为极限环, 仅为闭合曲线.\\

\subsubsection{\large 极限环的存在性}
给定一个方程组, 如何判断是否存在极限环?\\
先讨论如何判断一个方程组不存在极限环.\\
\textbf{本迪克松准则}\\
假设:D是平面的一个区域, 计算向量场散度
\[
div \vec{F} = f_x + g_y
\]
若$ div \vec{F} \neq  0 $ , 则区域D内不存在极限环, 甚至不存在闭合轨迹.\\
证明:\\
反证法: 假设存在闭合轨迹$C \subset D$
\[
\because \vec{F} \hat{n} \equiv 0 
\]
\begin{align*}
    \therefore 
    \oint _C \vec{F} \cdot \hat{n} ds & = 0 \\
    \iint _D div \vec{F} dA & > 0 \hspace*{0.5em}or < 0
\end{align*}
矛盾!\\[0.5em]
eg:
\[
\begin{cases}
    x' = x^2 + y^2 + 1 \\
    y' = x^2 - y^2 
\end{cases}
\]
\[
div \vec{F} = 2x - 2y 
\]
\begin{figure}[htbp]
    \centering
\includegraphics[width=0.3\textwidth]{C:/Users/37/Desktop/Notes on Common Calculus Equations/Notes and pictures/cwf13.2.1.png}
    \caption{极限环例子}
    \label{fig:cwf13.2.1}
\end{figure}

如图情况, 无法作结论.\\
\textbf{临界点准则}
\[
\begin{cases}
D \longrightarrow \text{xy平面} \\
C \longrightarrow  \text{闭合轨迹}  \hspace*{0.5em} C \subset D
\end{cases}
\Rightarrow \text{在C内至少存在一个临界点.}\\
\]
利用逆否命题:\\
如果D不存在临界点, 则不存在闭合轨迹, 故不存在极限环.\\
显然 $x' = x^2 + y^2 + 1$ 不存在零值\\
所以这例题不存在极限环.\\[0.5em]
下介绍一个未完全解决的问题:
\[
\begin{cases}
    x' = ax^2 + bxy + cy^2 + dx + ey + f \\
    y' = a'x^2 + b'xy + c'y^2 + d'x + e'y + f'
\end{cases}
\]
一个二次方程组, 其系数为实数, 最多能有多少个极限环?\\
目前已知的最大数量为4, 但尚未证明是否存在5个或更多的极限环.

\newpage
\subsection{\LARGE \hspace*{0.1em}非线性方程组和一阶方程的关系}
\subsubsection{\large 简单介绍}
考虑如下方程组:
\[
\begin{cases}
    \frac{dx}{dt} = f(x,y) \\
    \frac{dy}{dt} = g(x,y)
\end{cases}
\]
几何: 速度场
\[ \vec{F} = f\hat{i} + g\hat{j} 
\hspace*{1em}
\begin{cases}
    x = x(t) \\
    y = y(t)
\end{cases}
\]
消去t, 得:
\[
\frac{dy}{dx} = \frac{g(x,y)}{f(x,y)}
\]
为方向场.\\
期望: $y=f(x)$, 显式解;\\
现象: $h(x,y) = 0$, 隐式解.\\
eg:
\[
\begin{cases}
    x' = y \\
    y' = -x
\end{cases}
\]
解:
\[
\begin{bmatrix}
    x \\
    y
\end{bmatrix} = c_1
\begin{bmatrix}
    \cos t \\
    -\sin t
\end{bmatrix} + c_2
\begin{bmatrix}
    \sin t \\
    \cos t
\end{bmatrix}
\]
消去t, 得:
\[
\frac{dy}{dx} = \frac{-x}{y}
\] 
\[ x^2 + y^2 = c \]\\[0.5em]
eg:
\[
\begin{cases}
    x' = -ax + bxy \\
    y' = cy - dxy
\end{cases}
\]
1. 找临界点
\[
\begin{cases}
    x(-a + by) = 0 \\
    y(c - dx) = 0
\end{cases}
\]
解:
\[
\begin{cases}
    x = 0 \\
    y = 0 
\end{cases}
\hspace*{0.5em} or 
\begin{cases}
    x = \frac{c}{d} \\
    y = \frac{a}{b}
\end{cases}
\]
2. 线性化\\
对$(0,0)$线性化:
\[
\begin{bmatrix}
    -a & 0 \\
    0 & c
\end{bmatrix}
\hspace*{0.5em}
\lambda = -a, c
\]
鞍面\\
为简化计算, 设$a=b=c=d=1$
\[
\begin{cases}
    x' = -x + xy \\
    y' = y - xy
\end{cases}
\]
\[
\mathcal{J} =
\begin{bmatrix}
    -1 + y & x \\
    -y & 1 - x
\end{bmatrix}
\]
对$(\frac{c}{d},\frac{a}{b})$, 现$(1,1)$线性化:
\[
\mathcal{J}_{(1,1)} =
\begin{bmatrix}
    0 & 1 \\
    -1 & 0
\end{bmatrix}
\]
线性化的方程组是一个中心(一个个紧挨着的圆组成)\\[0.5em]
\textbf{沃尔泰拉法则}
\[
\begin{cases}
    x' = -x + xy \\
    y' = y - xy
\end{cases}
\]
3.消去$t$
\[
\frac{dy}{dx} = \frac{y(1-x)}{x(-1+y)}
\]
\begin{align*}
    \frac{y-1}{y} dy & = \frac{1-x}{x} dx \\
    y-\ln y & = \ln x - x + c_1 \\
    \frac{e^y}{y} & = \frac{x}{e^x} + c_2\\
    \frac{x}{e^x}\cdot \frac{y}{e^y} & = C
\end{align*}
可视为等值线:
\[
x e^{-x}\cdot y e^{-y} = h(x,y)
\]
最后:
\[
\begin{cases}
    x' = -ax + bxy - kx = -(a+k)x + bxy \\
    y' = cy - dxy - ky = (c-k)y - dxy
\end{cases}
\]
旧临界点$(\frac{c}{d},\frac{a}{b})$, 新临界点$(\frac{c-k}{d},\frac{a+k}{b})$\\
\begin{figure}[htbp]
    \centering
\includegraphics[width=0.5\textwidth]{C:/Users/37/Desktop/Notes on Common Calculus Equations/Notes and pictures/cwf13.2.2.png}
    \caption{临界点变化}
    \label{fig:cwf13.2.2}
\end{figure}

新的临界点显现了某一值降低的现象, 平衡偏差, 此现象称作沃尔泰拉法则.\\
至此本文主脉络已全部完成, 其中有意思的章节将额外补充.